%%%%%%%%%%%%%%%%%%%%%%%%%%%%%%%%%%%%%%%%%%%%%%%%%%%%%%%%%%%%%%%%%%%%%%%%%%%%%%%%%%
\begin{frame}[fragile]\frametitle{}

\begin{center}
{\Large STL Inspector App Workflow}
\end{center}
\end{frame}

%%%%%%%%%%%%%%%%%%%%%%%%%%%%%%%%%%%%%%%%%%%%%%%%%%%%%%%%%%%
%% WORKSHOP OVERVIEW
%%%%%%%%%%%%%%%%%%%%%%%%%%%%%%%%%%%%%%%%%%%%%%%%%%%%%%%%%%%
\begin{frame}[fragile]\frametitle{Use-case}
% \textbf{What we will build:}
A STEP/STL geometry inspector \& validator, \texttt{stlinspector}, with a CLI and a Streamlit app. Meshes exported from CAD/CAM tools can silently carry bad geometry -- non-watertight shells, non-manifold edges, degenerate faces -- that only surfaces later, as a failed 3D print or a broken CAM toolpath. \texttt{stlinspector} catches these issues right after export, with both a scriptable CLI (for CI pipelines) and a Streamlit app (for a quick visual check). This workshop builds it end-to-end in a single solo walkthrough, using Claude Code's structural primitives -- commands, subagents, skills, and hooks -- on top of the standard chat loop.

% \textbf{What we will learn:}
  % \begin{itemize}
    % \item Claude Code across 4 core SDLC stages: bootstrap, specification, architecture, implementation
    % \item Custom \textbf{commands}, \textbf{skills}, \textbf{subagents}, and \textbf{hooks}, beyond the standard chat loop
    % \item QA workflows (synthetic test fixtures)
    % \item Token-efficient habits, including ones that matter more on a larger, CAD/CAM-style codebase
  % \end{itemize}

% \textbf{Duration:} $\approx$120 min \quad \textbf{Format:} Solo walkthrough, one stage at a time

% \textbf{Stages covered:}
  % \begin{enumerate}
    % \item \textbf{Bootstrap}: \texttt{/init}, CLAUDE.md
    % \item \textbf{Specification}: subagent-written spec
    % \item \textbf{Architecture}: Plan mode, scaffold
    % \item \textbf{Implementation}: engine, CLI, Streamlit
    % \item \textbf{Commands \& Skills}: \lstinline|/add-check|, report-summary skill
    % \item \textbf{QA}: test subagent
    % \item \textbf{Hooks}: deterministic guardrails
    % \item \textbf{Best Practices \& Token Management}
  % \end{enumerate}
\end{frame}


%%%%%%%%%%%%%%%%%%%%%%%%%%%%%%%%%%%%%%%%%%%%%%%%%%%%%%%%%%%%%%%%%%%%%%%%%%%%%%%%%%
\begin{frame}[fragile]\frametitle{Claude Code Customization, at a Glance}
Four ways to extend Claude Code beyond the standard chat loop, used throughout this workshop:

\begin{center}
\adjustbox{max width=\linewidth}{%
\begin{tabular}{|l|l|l|}
\hline
\textbf{Primitive} & \textbf{Invoked via} & \textbf{Purpose} \\
\hline
Commands & Typed slash command, e.g. \texttt{/add-check} & Repeatable, team-shared workflow saved as a prompt \\
Subagents & Delegated by name, e.g. \texttt{specs}, \texttt{qa} & Isolate a task's context; only its summary returns \\
Skills & Auto-loaded when the task matches its description & Lazy-loaded domain knowledge, not in every session \\
Hooks & Fire automatically on lifecycle events & Deterministic guardrail -- always runs, not a request \\
\hline
\end{tabular}
}
\end{center}
\end{frame}


%%%%%%%%%%%%%%%%%%%%%%%%%%%%%%%%%%%%%%%%%%%%%%%%%%%%%%%%%%%
%% STAGE 0: SETUP
%%%%%%%%%%%%%%%%%%%%%%%%%%%%%%%%%%%%%%%%%%%%%%%%%%%%%%%%%%%
\begin{frame}[fragile]\frametitle{}
\begin{center}
{\Large Setup}
\end{center}
\end{frame}

%%%%%%%%%%%%%%%%%%%%%%%%%%%%%%%%%%%%%%%%%%%%%%%%%%%%%%%%%%%%%%%%%%%%%%%%%%%%%%%%%%
\begin{frame}[fragile]\frametitle{Pre-requisites}
% Before starting, verify all of the following:

\textbf{Essential Steps/Checklist:}
  \begin{itemize}
    \item Create an \textbf{empty folder} and switch to it
	\item \lstinline|mkdir stlinspector && cd stlinspector|
    \item Conda installed: \lstinline|conda --version|
    \item Claude Code installed: \lstinline|claude doctor|
    % \item Git configured: \lstinline|git config --global user.name "Your Name"|
    % \item Python $\geq$ 3.11: \lstinline|python --version|
    \item Have Claude Pro/Max or API key
	\item Create the conda env and install demo deps
  \end{itemize}

\begin{lstlisting}
conda create -n stlinspector python=3.11 -y
conda activate stlinspector

pip install trimesh numpy streamlit plotly pytest rich
\end{lstlisting}
\end{frame}


%%%%%%%%%%%%%%%%%%%%%%%%%%%%%%%%%%%%%%%%%%%%%%%%%%%%%%%%%%%%%%%%%%%%%%%%%%%%%%%%%%
\begin{frame}[fragile]\frametitle{Choosing the Right Model}
Claude Code ships three model tiers. Use the right one:

\begin{center}
\adjustbox{max width=\linewidth}{%
\begin{tabular}{|l|l|l|}
\hline
\textbf{Model} & \textbf{Best For} & \textbf{Avoid For} \\
\hline
Haiku 4.5 & commands, exploration, quick edits & Deep reasoning \\
Sonnet 4.6 (default) & Daily coding, refactors, implementation & Trivial lookups \\
Opus 4.6+ & Architecture, hard trade-off calls & Repetitive work \\
\hline
\end{tabular}
}
\end{center}

\textbf{Setting model}:
\begin{lstlisting}
/model claude-opus-4-6 # inside a session
claude --model claude-opus-4-6 # Or launch

# Per-project default in .claude/settings.json:
{ "model": "claude-sonnet-4-6" }
\end{lstlisting}
\end{frame}


%%%%%%%%%%%%%%%%%%%%%%%%%%%%%%%%%%%%%%%%%%%%%%%%%%%%%%%%%%%
%% PROJECT BOOTSTRAP
%%%%%%%%%%%%%%%%%%%%%%%%%%%%%%%%%%%%%%%%%%%%%%%%%%%%%%%%%%%
\begin{frame}[fragile]\frametitle{}
\begin{center}
{\Large Stage 1: Project Bootstrap}
\end{center}
\end{frame}


%%%%%%%%%%%%%%%%%%%%%%%%%%%%%%%%%%%%%%%%%%%%%%%%%%%%%%%%%%%%%%%%%%%%%%%%%%%%%%%%%%
\begin{frame}[fragile]\frametitle{Run /init: Generate CLAUDE.md}
  \begin{itemize}
    % \item Make sure \textbf{Plan mode} is active first (\texttt{Shift+Tab}). Project folder is empty except \texttt{.git}.
	\item Start \lstinline|claude| and at the prompt say \lstinline|/init| as shown below
    \item Claude asks clarifying questions about the project: answer them, then approve the file creation.
  \end{itemize}

\begin{lstlisting}
/init

## What /init generates:
CLAUDE.md              # Project intelligence file
.claude/settings.json  # Permissions, model config

## Resulting CLAUDE.md (edit to match below):
---
# CLAUDE.md

## Project: stlinspector
STL/STEP geometry inspector and validator (PoC).

## Tech Stack
- Python 3.11, trimesh, numpy
- Streamlit + plotly (3D preview), Rich (CLI)

## Build Commands
- Install: `pip install trimesh numpy streamlit plotly pytest rich`
- Test:    `pytest tests/ -v`
---
## Commit immediately:
git add .
git commit -m "chore: added CLAUDE.md"
\end{lstlisting}
\end{frame}


%%%%%%%%%%%%%%%%%%%%%%%%%%%%%%%%%%%%%%%%%%%%%%%%%%%%%%%%%%%
%% SPECIFICATIONS
%%%%%%%%%%%%%%%%%%%%%%%%%%%%%%%%%%%%%%%%%%%%%%%%%%%%%%%%%%%
\begin{frame}[fragile]\frametitle{}
\begin{center}
{\Large Stage 2: Specifications}
\end{center}
\end{frame}


%%%%%%%%%%%%%%%%%%%%%%%%%%%%%%%%%%%%%%%%%%%%%%%%%%%%%%%%%%%%%%%%%%%%%%%%%%%%%%%%%%
\begin{frame}[fragile]\frametitle{Specs Subagent Writes the Project Spec}

  \begin{itemize}
	\item Crate a \textbf{specs} subagent that keeps spec-writing out of your main context.
    \item \lstinline|mkdir -p .claude/agents|
    \item \lstinline|cat > .claude/agents/specs.md| as below
  \end{itemize}

\begin{lstlisting}
---
name: specs
description: >
  Technical specification writer. Produces structured Markdown specs from feature descriptions. 
Read-only.
model: claude-opus-4-6
allowed-tools: Read, Write, Glob

Produce a Markdown spec covering: Overview, Goals,
Non-Goals, Functional Requirements, Data Models,
API/CLI Contracts, Open Questions. 

Save to docs/specs/<feature>.md
---
## Invoke it via following prompt:
---
Write a technical spec for stlinspector. 
Cover:
- core.py: load_mesh(path), InspectionReport dataclass
  (bbox, volume, surface_area, triangle_count, issues)
- Validation rules: not-watertight, non-manifold edges,
  degenerate faces
- cli.py: python cli.py <file> [--report out.json]
  (no console-script entry point in this pass)
- app.py: upload -> 3D preview -> report table
Save to docs/specs/stlinspector-spec.md
---
git add docs/
git commit -m "docs: added stlinspector spec"
\end{lstlisting}
\end{frame}


%%%%%%%%%%%%%%%%%%%%%%%%%%%%%%%%%%%%%%%%%%%%%%%%%%%%%%%%%%%
%% ARCHITECTURE
%%%%%%%%%%%%%%%%%%%%%%%%%%%%%%%%%%%%%%%%%%%%%%%%%%%%%%%%%%%
\begin{frame}[fragile]\frametitle{}
\begin{center}
{\Large Stage 3: Architecture Design}
\end{center}
\end{frame}


%%%%%%%%%%%%%%%%%%%%%%%%%%%%%%%%%%%%%%%%%%%%%%%%%%%%%%%%%%%%%%%%%%%%%%%%%%%%%%%%%%
\begin{frame}[fragile]\frametitle{Design, Then Scaffold}
Plan mode is \textbf{read-only}: no edits, no bash, and is perfect for exploring the module split before committing to it.
\lstinline|Shift+Tab   # Enter Plan mode|

\begin{lstlisting}
-- Prompt --
Read docs/specs/stlinspector-spec.md, then design the package layout for stlinspector:
- src/{core,cli,app,config}.py
- Data flow: file -> core.load_mesh() -> mesh -> core.inspect_mesh() -> InspectionReport
- What core.py must NOT depend on (no Streamlit imports)
Show the plan before creating any files.
--
## Approve ("Yes, auto-accept edits")

\lstinline|Shift+Tab   # Enter Normal mode|

-- Prompt --
Create the stlinspector scaffold with stub methods:
- src/core.py    (stub: load_mesh(path), inspect_mesh(mesh))
- src/cli.py     (stub: main())
- src/app.py     (stub: Streamlit skeleton)
- src/config.py  (stub: Config dataclass)
- tests/__init__.py, tests/conftest.py, tests/test_core.py
- README.md

---

git add . && git commit -m "chore: scaffold stlinspector"
/cost      # Check token usage so far
\end{lstlisting}
\end{frame}


%%%%%%%%%%%%%%%%%%%%%%%%%%%%%%%%%%%%%%%%%%%%%%%%%%%%%%%%%%%
%% IMPLEMENTATION
%%%%%%%%%%%%%%%%%%%%%%%%%%%%%%%%%%%%%%%%%%%%%%%%%%%%%%%%%%%
\begin{frame}[fragile]\frametitle{}
\begin{center}
{\Large Stage 4: Implementation}
\end{center}
\end{frame}


%%%%%%%%%%%%%%%%%%%%%%%%%%%%%%%%%%%%%%%%%%%%%%%%%%%%%%%%%%%%%%%%%%%%%%%%%%%%%%%%%%
\begin{frame}[fragile]\frametitle{Implement core.py: the Inspection Engine}
\begin{lstlisting}
@file:src/core.py
@file:docs/specs/stlinspector-spec.md
Implement core.py:
- load_mesh(path: str) -> trimesh.Trimesh
- inspect_mesh(mesh: trimesh.Trimesh) -> InspectionReport, where InspectionReport is a dataclass with: bounding_box, volume, surface_area, triangle_count, issues: list[str]
- Issue checks:
  - "not_watertight" if not mesh.is_watertight
  - "non_manifold_edges" if any edge shared by > 2 faces
  - "degenerate_faces" if any face area < 1e-8
- Full type hints, Google-style docstrings
- Raise FileNotFoundError / ValueError with clear messages

---
git add src/core.py
git commit -m "feat: implement inspection engine"
\end{lstlisting}
\end{frame}


%%%%%%%%%%%%%%%%%%%%%%%%%%%%%%%%%%%%%%%%%%%%%%%%%%%%%%%%%%%%%%%%%%%%%%%%%%%%%%%%%%
\begin{frame}[fragile]\frametitle{Implement cli.py and app.py}
\begin{lstlisting}
## CLI:
@file:src/cli.py
Implement the CLI (run via `python cli.py`, no
console-script entry point in this pass):
- load_mesh(path) then inspect_mesh(mesh) [--report out.json] via argparse
- Use Rich Table to print bbox/volume/area/triangle_count
- Print issues in red if any, green "All checks passed"
- Exit code 1 if any issue found (CI-friendly)
- JSON is the only report format (via --report out.json)
  -- this feeds the inspection-report-summary skill later

---
## Streamlit:
@file:src/app.py
Implement the Streamlit app:
- st.file_uploader restricted to .stl
- Save upload to a temp file, call core.load_mesh() then core.inspect_mesh()
- 3D preview via plotly go.Mesh3d (mesh.vertices/faces)
- st.metric tiles for bbox/volume/area/triangle_count
- st.error per issue, st.success if none

---
## Run and verify (no entry point installed yet):
python src/cli.py tests/fixtures/box.stl --report out.json
streamlit run src/app.py

---
git add src/ && git commit -m "feat: implement CLI and Streamlit app"
\end{lstlisting}
\end{frame}


%%%%%%%%%%%%%%%%%%%%%%%%%%%%%%%%%%%%%%%%%%%%%%%%%%%%%%%%%%%
%% COMMANDS AND SKILLS
%%%%%%%%%%%%%%%%%%%%%%%%%%%%%%%%%%%%%%%%%%%%%%%%%%%%%%%%%%%
\begin{frame}[fragile]\frametitle{}
\begin{center}
{\Large Custom Commands \& Skills}
\end{center}
\end{frame}


%%%%%%%%%%%%%%%%%%%%%%%%%%%%%%%%%%%%%%%%%%%%%%%%%%%%%%%%%%%%%%%%%%%%%%%%%%%%%%%%%%
\begin{frame}[fragile]\frametitle{Create the /add-check Custom Command}
  \begin{itemize}
	\item If you would repeat the same ``add a validation rule'' prompt more than once, make it a command.
	\item \lstinline|mkdir -p .claude/commands|
	\item Write \lstinline|.claude/commands/add-check.md| as below
  \end{itemize}

\begin{lstlisting}
---
description: >
  Scaffold a new geometry validation rule in core.py,
  wired into InspectionReport.issues, plus a matching
  synthetic-fixture test. Pass the rule name as argument.
allowed-tools: Read, Edit, Write, Bash

Add a new validation check named "$ARGUMENTS" to
src/core.py:
1. Add a private _check_$ARGUMENTS(mesh) -> bool helper
2. Wire it into inspect(), appending "$ARGUMENTS" to
   issues when it returns True
3. Add a synthetic-mesh test in tests/test_core.py that
   triggers the new check
4. Run pytest tests/ -v and report the result
---

## Usage:
/add-check thin_walls
\end{lstlisting}
\end{frame}


%%%%%%%%%%%%%%%%%%%%%%%%%%%%%%%%%%%%%%%%%%%%%%%%%%%%%%%%%%%%%%%%%%%%%%%%%%%%%%%%%%
\begin{frame}[fragile]\frametitle{Create the inspection-report-summary Skill}

  \begin{itemize}
	\item \textbf{Skills} are lazy-loaded: Claude pulls this in only when the task matches, unlike CLAUDE.md which loads every session.
	\item A skill cannot do geometry math itself -- \texttt{core.py}/trimesh already did that. What it \textit{can} do well is turn a report into a stakeholder-ready summary.
	\item \lstinline|mkdir -p .claude/skills/inspection-report-summary|
	\item Write \lstinline|SKILL.md| in that folder, as below
  \end{itemize}


\begin{lstlisting}
---
name: inspection-report-summary
description: >
  Summarizes an stlinspector InspectionReport (JSON)
  for a non-technical stakeholder. Use when asked to
  "summarize", "explain", or "triage" an inspection
  report.
---

## Input
Read the JSON report written by --report (e.g. out.json,
or the path given): bbox, volume, surface_area,
triangle_count, and a list of issues (if any).

## Output
- One-paragraph plain-English verdict: PASS / FAIL and why
- If issues present, rank by severity:
  1. not_watertight (blocks 3D printing / CAM)
  2. non_manifold_edges (blocks CAM toolpath generation)
  3. degenerate_faces (usually cosmetic, low risk)
- Suggest the next concrete action per issue found
  (e.g. "re-export from CAD with tighter tessellation")

## Generate a report to summarize (if out.json isn't already present):
python src/cli.py src/Box.stl --report out.json

## Usage:
Summarize out.json for a non-technical reviewer
\end{lstlisting}
\end{frame}


%%%%%%%%%%%%%%%%%%%%%%%%%%%%%%%%%%%%%%%%%%%%%%%%%%%%%%%%%%%
%% QA WORKFLOWS
%%%%%%%%%%%%%%%%%%%%%%%%%%%%%%%%%%%%%%%%%%%%%%%%%%%%%%%%%%%
\begin{frame}[fragile]\frametitle{}
\begin{center}
{\Large QA Workflows}
\end{center}
\end{frame}


%%%%%%%%%%%%%%%%%%%%%%%%%%%%%%%%%%%%%%%%%%%%%%%%%%%%%%%%%%%%%%%%%%%%%%%%%%%%%%%%%%
\begin{frame}[fragile]\frametitle{QA Subagent}

  \begin{itemize}
	\item Synthetic Fixtures, No Sample Files
	\item Write \lstinline|.claude/agents/qa.md| as below.
	\item Use the qa subagent to write tests/conftest.py with fixtures: \lstinline|valid_box_mesh, non_watertight_mesh, degenerate_face_mesh|. Then write \lstinline|tests/test_core.py| covering all three cases against \lstinline|core.inspect()|.
  \end{itemize}
 
 
\begin{lstlisting}
---
name: qa
description: >
  Generates pytest tests using synthetic trimesh fixtures
  (no real STL files). Use when adding tests for core.py.
model: claude-sonnet-4-6
allowed-tools: Read, Write, Bash

Generate fixtures with trimesh.creation (box, icosphere)
for the valid case, and hand-built meshes (remove a face for
non-watertight, zero-area triangle for degenerate) for
the failing cases. Assert on InspectionReport fields
and on report.issues membership.
---
\end{lstlisting}
\end{frame}


%%%%%%%%%%%%%%%%%%%%%%%%%%%%%%%%%%%%%%%%%%%%%%%%%%%%%%%%%%%%%%%%%%%%%%%%%%%%%%%%%%
\begin{frame}[fragile]\frametitle{Run Tests and Check Coverage}
\begin{lstlisting}[language=bash]
pytest tests/ -v --tb=short

pytest tests/ --cov=src --cov-report=term-missing

--- 
# If a test fails, hand Claude the traceback directly:
---
@file:src/core.py
test_degenerate_face_mesh fails with:
AssertionError: 'degenerate_faces' not in report.issues
Here is the fixture: [paste]. Fix the check or the fixture.
---
git add tests/
git commit -m "test: add synthetic-fixture test suite"
\end{lstlisting}
\end{frame}




%%%%%%%%%%%%%%%%%%%%%%%%%%%%%%%%%%%%%%%%%%%%%%%%%%%%%%%%%%%
%% HOOKS
%%%%%%%%%%%%%%%%%%%%%%%%%%%%%%%%%%%%%%%%%%%%%%%%%%%%%%%%%%%
\begin{frame}[fragile]\frametitle{}
\begin{center}
{\Large Hooks: Deterministic Guardrails}
\end{center}
\end{frame}


%%%%%%%%%%%%%%%%%%%%%%%%%%%%%%%%%%%%%%%%%%%%%%%%%%%%%%%%%%%%%%%%%%%%%%%%%%%%%%%%%%
\begin{frame}[fragile]\frametitle{Wire Up Hooks in .claude/settings.json}
  \begin{itemize}
	\item Hooks are deterministic: they always fire, unlike a prompt-based instruction.
	\item What this buys:
	  \begin{itemize}
		\item PostToolUse: every edited file is auto-formatted
		\item Stop: the suite must pass before Claude can finish
		\item PreToolUse: destructive bash commands are blocked even if Claude "forgets", no prompt can override this
	  \end{itemize}
  \end{itemize}


\begin{lstlisting}
{
  "hooks": {
    "PostToolUse": [{
      "matcher": "Edit|Write",
      "hooks": [{"type": "command",
        "command": "FILE=$(cat | jq -r '.tool_input.file_path'); ruff format \"$FILE\" 2>/dev/null || true"}]
    }],
    "Stop": [{
      "hooks": [{"type": "command",
        "command": "cd \"$CLAUDE_PROJECT_DIR\" && pytest tests/ -q --tb=short || echo '{\"decision\": \"block\", \"reason\": \"tests failing\"}'"}]
    }],
    "PreToolUse": [{
      "matcher": "Bash",
      "hooks": [{"type": "command",
        "command": "INPUT=$(cat); CMD=$(echo $INPUT | jq -r '.tool_input.command'); echo \"$CMD\" | grep -qE 'rm -rf|DROP TABLE' && echo 'Blocked: destructive command' >&2 && exit 2 || exit 0"}]
    }]
  }
}
\end{lstlisting}
\end{frame}


%%%%%%%%%%%%%%%%%%%%%%%%%%%%%%%%%%%%%%%%%%%%%%%%%%%%%%%%%%%
%% DEMO WRAP-UP
%%%%%%%%%%%%%%%%%%%%%%%%%%%%%%%%%%%%%%%%%%%%%%%%%%%%%%%%%%%
\begin{frame}[fragile]\frametitle{Run End-to-End and Recap}
\textbf{What we covered:} bootstrap $\to$ spec $\to$ architecture $\to$ implementation $\to$ commands \& skills $\to$ QA subagent $\to$ hooks -- everything needed to ship \texttt{stlinspector} end-to-end.

\begin{lstlisting}[language=bash]
python src/cli.py tests/fixtures/box.stl --report out.json
# Expected: Rich table with bbox/volume/area/triangle_count
# "All checks passed" in green (box is watertight)
# Writes out.json (the only report format)

streamlit run src/app.py
# Expected: browser opens at http://localhost:8501
# Upload an STL, see 3D preview + report

pytest tests/ -v      # All synthetic-fixture tests pass
\end{lstlisting}
\end{frame}


%%%%%%%%%%%%%%%%%%%%%%%%%%%%%%%%%%%%%%%%%%%%%%%%%%%%%%%%%%%
%% BEST PRACTICES AND TOKEN MANAGEMENT (SELF-CONTAINED)
%%%%%%%%%%%%%%%%%%%%%%%%%%%%%%%%%%%%%%%%%%%%%%%%%%%%%%%%%%%
\begin{frame}[fragile]\frametitle{}
\begin{center}
{\Large Best Practices \& Token Management}
\end{center}
\end{frame}


%%%%%%%%%%%%%%%%%%%%%%%%%%%%%%%%%%%%%%%%%%%%%%%%%%%%%%%%%%%%%%%%%%%%%%%%%%%%%%%%%%
\begin{frame}[fragile]\frametitle{Token Budgeting and Model Choice}
  \begin{itemize}
    \item Every input \textit{and} output token is billed: unfocused sessions on Sonnet/Opus cost more than necessary.
    \item Bloated context buries signal under failed attempts and irrelevant files; a smaller, focused context produces sharper output.
    \item \textbf{Rule of thumb:} default to Sonnet; drop to Haiku for subagents doing \texttt{qa} triage runs); promote to Opus only for real architecture trade-offs.
  \end{itemize}
\begin{lstlisting}[language=bash]
/model claude-haiku-4-5           # Cheap subagent work
/cost                              # Check spend so far
\end{lstlisting}
\end{frame}


%%%%%%%%%%%%%%%%%%%%%%%%%%%%%%%%%%%%%%%%%%%%%%%%%%%%%%%%%%%%%%%%%%%%%%%%%%%%%%%%%%
\begin{frame}[fragile]\frametitle{Plan Mode and Breaking Tasks Down}
\textbf{Plan mode} (\texttt{Shift+Tab}) is read-only: exploring a wrong approach costs almost nothing; implementing it wrong costs a redo. Stay in \textbf{normal mode} (not auto-accept) for multi-step work; only auto-accept mechanical edits already reviewed in Plan mode.

\begin{lstlisting}
I want to add a wall-thickness check to stlinspector.
Break this into numbered subtasks and show me the list
first. Do NOT start until I confirm each step.

## TodoWrite tracks progress automatically as steps
## complete, one step in progress at a time, no
## parallel context drift.
\end{lstlisting}
\end{frame}


%%%%%%%%%%%%%%%%%%%%%%%%%%%%%%%%%%%%%%%%%%%%%%%%%%%%%%%%%%%%%%%%%%%%%%%%%%%%%%%%%%
\begin{frame}[fragile]\frametitle{Context Hygiene}
  \begin{itemize}
    \item \texttt{/compact}: summarise older turns once \texttt{/cost} shows context past $\sim$50\% of the window.
    \item \texttt{@file:} scoping: load only the file(s) the current step needs: loading the whole \texttt{src/} tree for a one-file fix is the most common waste.
    \item Subagents (\texttt{specs}, \texttt{qa}) run in isolated contexts: only their summary returns to the main session.
    \item Skills are lazy-loaded: \texttt{inspection-report-summary} only enters context when the task matches.
    \item \texttt{/clear}: nuclear option, only for a genuinely unrelated new task.
  \end{itemize}
\begin{lstlisting}[language=bash]
@file:src/core.py   # Not the whole src/ tree
/compact
\end{lstlisting}
\end{frame}


%%%%%%%%%%%%%%%%%%%%%%%%%%%%%%%%%%%%%%%%%%%%%%%%%%%%%%%%%%%%%%%%%%%%%%%%%%%%%%%%%%
\begin{frame}[fragile]\frametitle{Structure Beats Prompting Skill}
Everything used in this workshop is a structural layer, not a clever prompt:
  \begin{itemize}
    \item \texttt{CLAUDE.md}: execution memory, loaded every session (WHY/WHAT/HOW, kept short)
    \item \texttt{.claude/commands/}: repeated workflows as \texttt{/add-check}-style commands, committed to Git for the whole team
    \item \texttt{.claude/rules/}: path-scoped constraints near sharp edges (e.g. \texttt{core.py} numerics)
    \item \texttt{.claude/skills/}: lazy-loaded domain knowledge (\texttt{inspection-report-summary})
    \item Hooks: deterministic guardrails, not requests
  \end{itemize}
\textbf{Prompts are temporary. Structure is permanent}: it is what let Stages 5 to 8 of this workshop run with almost no re-explaining.
\end{frame}


%%%%%%%%%%%%%%%%%%%%%%%%%%%%%%%%%%%%%%%%%%%%%%%%%%%%%%%%%%%%%%%%%%%%%%%%%%%%%%%%%%
\begin{frame}[fragile]\frametitle{Session Continuity Across Days}
A workshop like this one often spans multiple sittings: treat sessions as resumable, not disposable.
\begin{lstlisting}[language=bash]
/rename stlinspector-implementation   # Name current session
claude --resume stlinspector-implementation
claude -c                             # Resume most recent

## Or: write progress to a plan file Claude can re-read:
Write our progress to .claude/plans/current.md as a
numbered checklist with [x]/[ ] and the next step noted
at the top.
\end{lstlisting}
\end{frame}


%%%%%%%%%%%%%%%%%%%%%%%%%%%%%%%%%%%%%%%%%%%%%%%%%%%%%%%%%%%%%%%%%%%%%%%%%%%%%%%%%%
\begin{frame}[fragile]\frametitle{Scaling Up: Large CAD/CAM Codebases}
\texttt{stlinspector} is a small PoC. Real CAD/CAM systems (geometry kernels, STEP/IGES parsers, GUI layers) can be millions of lines, so the same habits matter more, not less:
  \begin{itemize}
    \item \textbf{Module-scoped subagents:} one for the geometry kernel, one for file-format parsers, one for the GUI, never one subagent with the whole repo in scope.
    \item \textbf{Directory-scoped rules files:} e.g.\ \texttt{.claude/rules/\allowbreak kernel-rules.md} with \texttt{paths: ["kernel/**"]}, so numerics constraints load only when kernel files are touched.
    \item \textbf{\texttt{@file:} precision:} never ask Claude to ``read the parser'', instead reference the exact translation-unit or module; a full STEP/IGES parser can be tens of thousands of lines.
  \end{itemize}
\end{frame}


%%%%%%%%%%%%%%%%%%%%%%%%%%%%%%%%%%%%%%%%%%%%%%%%%%%%%%%%%%%%%%%%%%%%%%%%%%%%%%%%%%
\begin{frame}[fragile]\frametitle{Scaling Up: Output Token Discipline}
  \begin{itemize}
    \item \textbf{Ask for diffs, not rewrites:} ``show the changed lines'' beats ``rewrite the file'': kernel files are often thousands of lines; a full rewrite burns output tokens and risks silently dropping code.
    \item \textbf{Synthetic fixtures over real files:} as in this workshop, use \texttt{trimesh.creation} primitives, not multi-MB STEP files, as test/context input.
    \item \textbf{Plan mode is mandatory before touching kernel/numerics code}: a wrong edit to tolerance-sensitive geometry code is expensive to detect and worse to revert.
    \item \textbf{One validation rule (or one kernel change) per request}: mirrors the \texttt{/add-check} pattern; batching multiple kernel changes in one prompt compounds risk and review cost.
  \end{itemize}
\end{frame}


